\section{The Documented Patient Concerns}

\subsection{Where the twenty-one come from}

\draftinstr{Roughly 1,000 characters. Two dated research passes on July 28,
2026, one by Gemini 3.1 Pro and one by ChatGPT 5.6 Thinking Extended, filed at
\texttt{research/research-a.md} and \texttt{research/research-b.md}. Six concern
families and sixteen numbered concerns, overlapping on five items,
de-duplicating to twenty-one. Cite the underlying human evidence:
\texttt{WuSemiAI2026}, \texttt{TelesAI2026A}, \texttt{Jauniaux2025}, and
\texttt{BrarRAS2024X}. State the two headline numbers: 56 percent of 50
previously operated patients would consider semi-autonomous robot-assisted
surgery, conditional on which steps the robot controls; and approximately half
of 330 oncology patients expressed concern about artificial intelligence in
cancer care.}

\vspace{0.30em}

\noindent\begin{tabularx}{\textwidth}{L{4.3cm} C{1.5cm} Y}
\toprule
\textbf{Concern family} & \textbf{Concerns} & \textbf{Where this paper answers it} \\
\midrule
Control and the human element & 5 & \S7.3, \S7.4; Figures 8, 18, 19 \\
Safety and technical failure & 5 & \S7.4, \S7.5, \S7.6; Figures 20, 21 \\
Evidence and effectiveness & 4 & \S4.2, \S5.1, \S10; Figures 13, 27, 28 \\
Data, privacy, and security & 3 & \S9.3, \S9.4; Figures 17, 24 \\
Fairness and applicability & 2 & \S6.3, \S8.3; Figures 14, 23 \\
Accountability and burden & 2 & \S11, \S12; Figures 29, 30 \\
\bottomrule
\end{tabularx}

\vspace{0.30em}

\draftinstr{Verify the concern count sums to 21 and that every family row names
at least one section and one figure. Source the mapping from
\texttt{research/README.md}, which carries both concern-to-section tables.}

\begin{pafig}[d2-type: nested containers]
\node[d2key] (a) at (0,0) {Figure 5 slot\\six concern containers};
\node[d2box] (b) at (0,-1.8) {Source:\\d2/fig-05-concern-families.d2};
\draw[d2edge] (a) -- (b);
\end{pafig}
\vspace{-0.7cm}
\figcaption{Figure 5. The twenty-one documented concerns grouped into six disjoint
families, with the answering section and figure printed inside each item, and a fill\\
key stating how many are answered by a limit, a guarantee, or governance alone.}

\draftinstr{Replace the Figure 5 slot with the full D2-type nested-container
figure from \texttt{d2/fig-05-concern-families.d2}. Six \texttt{d2cont}
containers in a 3 by 2 arrangement at the coordinates given in that file's TikZ
notes. Fill key: seven items Corporate Blue for a hard numeric limit, nine white
for a procedural guarantee, five \texttt{pagraym} for governance or disclosure
only.}

\subsection{The five families, one subsection each}

\subsubsection{Control and the human element}

\draftinstr{Roughly 1,400 characters. Concerns 1 to 5: who is driving, can the
surgeon stop it, will the surgeon defer to the model, will I still have a doctor
who knows me, less time with my care team. Cite \texttt{WuSemiAI2026} and
\texttt{BrarRAS2024X}; the latter reports that the most common public
misconception is that the robot operates by itself. Answer with the Class II
collaborative classification and the absence of a model-to-actuator path, and
forward to Figures 8 and 19.}

\subsubsection{Safety and technical failure}

\draftinstr{Roughly 1,400 characters. Concerns 6 to 10: malfunction or
unintended motion, misreading tissue or a margin, injury to a major vessel,
instrument failure mid-procedure, conversion to open surgery. Answer with the
force caps, the vascular no-fly corridor, and the fault tree, and state plainly
that conversion to open surgery is a recorded outcome rather than a failure.
Cite \texttt{iec8060127} and \texttt{ieee7009} for the standards the safety
architecture is built against.}

\subsubsection{Evidence and effectiveness}

\draftinstr{Roughly 1,400 characters. Concerns 11 to 14: cancer control,
being among the first, newer described as better, and team experience. Cite
\texttt{DutchCohort2025} for the learning-curve comparator, \texttt{Jauniaux2025}
for the systematic divergence between expectation and experience, and
\texttt{FDARobot2022} for the FDA's own caution. This is the family the paper
answers least completely, and \S3.4 must say so.}

\subsubsection{Data, privacy, and security}

\draftinstr{Roughly 1,200 characters. Concerns 15 to 17: who sees the recording,
will my data train other systems, could the system be attacked. Answer with the
on-premises boundary, the per-store read lists, and the retention periods. Cite
\texttt{FDACyber2026} and \texttt{FDATrans2024}, and forward to Figures 17
and 24.}

\subsubsection{Fairness, applicability, accountability, and burden}

\draftinstr{Roughly 1,200 characters. Concerns 18 to 21: tested on people like
me, software change after consent, who is accountable, and what it costs. Cite
\texttt{FDADiver2017} for subgroup reporting and \texttt{FDARealW2025} for
real-world performance monitoring. Forward to Figures 14, 23, 29, and 30.}

\begin{pafig}[graphviz-type: bipartite graph, full page]
\node[gvkey] (a) at (0,0) {Figure 6 slot\\concern to clause map};
\node[gvbox] (b) at (0,-2.0) {Source:\\graphviz/fig-06-concern-to-clause.dot};
\draw[gvedge] (a) -- (b);
\end{pafig}
\vspace{-0.7cm}
\figcaption{Figure 6. Every documented concern wired to the clause, limit, or number that
answers it, drawn as a strictly bipartite graph so that no concern is answered with\\
another concern, and dashed where the answer is governance or disclosure alone.}

\draftinstr{Replace the Figure 6 slot with the full-page graphviz-type bipartite
map from \texttt{graphviz/fig-06-concern-to-clause.dot}. Twenty-one concern
ellipses on the left, nineteen clause boxes on the right, twenty-four edges
between them and none within a partition. Fourteen edges solid Corporate Blue,
seven solid gray, three dashed. This is the load-bearing figure of the paper.}

\subsection{How completely each concern is answered}

\draftinstr{Roughly 900 characters introducing Figure 7 and its two axes:
prevalence in the surveyed sources, and completeness of the protocol answer
scored 1.0 for a numeric limit, 0.7 for a procedural guarantee, 0.5 for a
governance commitment, and 0.3 for disclosure only. Name the four concerns in
quadrant 4 and do not argue them away.}

\begin{pafig}[mermaid-type: quadrantChart]
\node[mmgoal] (a) at (0,0) {Figure 7 slot\\concern quadrant};
\node[mmin] (b) at (0,-2.0) {Source:\\mermaid/fig-07-concern-quadrant.md};
\draw[mmedge] (a) -- (b);
\end{pafig}
\vspace{-0.7cm}
\figcaption{Figure 7. Sixteen documented concerns positioned by how often they are
raised in the surveyed literature against how completely this protocol answers them,\\
with the lower-right quadrant collecting the concerns answered by governance alone.}

\draftinstr{Replace the Figure 7 slot with the mermaid-type quadrant from
\texttt{mermaid/fig-07-concern-quadrant.md}. Fill the four quadrant fields
first, then the dividing cross, then the sixteen points as 1.6 mm Corporate Blue
discs. Where two points fall within 8 mm, alternate the label anchor and add a
0.4pt leader so no label overlaps another.}

\subsection{The same safeguards, read as capabilities you hold}

\draftinstr{Roughly 800 characters introducing Figure 8. Every safeguard in the
parent protocol is written with the sponsor as the actor. Reassigning the
primary actor is the intervention. Ten of the twelve capabilities are exercised
by the participant; the two that are not require someone physically present or
inside the build system, and neither can be exercised against the participant.}

\begin{pafig}[plantuml-type: use case, @startuml]
\node[umlkey] (a) at (0,0) {Figure 8 slot\\patient as primary actor};
\node[umlbox] (b) at (0,-1.9) {Source:\\plantuml/fig-08-patient-actor-usecase.puml};
\draw[umlarrow] (a) -- (b);
\end{pafig}
\vspace{-0.7cm}
\figcaption{Figure 8. The twelve safeguards of the protocol redrawn as a use case diagram
with the participant as the primary actor, so that each one reads as a capability they\\
hold rather than as an obligation the sponsor has undertaken on their behalf.}

\draftinstr{Replace the Figure 8 slot with the PlantUML-type use case diagram
from \texttt{plantuml/fig-08-patient-actor-usecase.puml}. Participant actor at
$(-6.4,-4.2)$, left of the boundary; three secondary actors right of it. Twelve
\texttt{umlusecase} ellipses in three columns and four rows. The two long
associations to the far column route with looseness 0.85 so they pass between
rows rather than through an ellipse.}

\subsection{Where each concern physically lives}

\draftinstr{Roughly 700 characters introducing Figure 9. A concern with no
physical address is unanswerable; a concern that resolves to a named cabinet, a
named console, or a named cable can be inspected. Source from
\texttt{diagrams-python/fig\_09\_concern\_locations.py} and its
\texttt{CONCERN\_ADDRESS} mapping.}

\begin{pafig}[diagrams (python)-type: Cluster + Node]
\node[dgtiled] (a) at (0,0) {};
\glyphai{a}{protowhite}
\node[dglabel,anchor=north] at ([yshift=-5.4mm]a.center) {Figure 9 slot, concern addresses};
\node[dglabel,anchor=north] at (0,-1.6) {Source: diagrams-python/fig\_09\_concern\_locations.py};
\end{pafig}
\vspace{-0.7cm}
\figcaption{Figure 9. Each documented concern resolved to the physical component that
answers it, grouped into the operating room, the on-premises cabinet, the governance\\
bodies, and everything outside the building, with the intra-procedural route absent.}

\draftinstr{Replace the Figure 9 slot with the full Diagrams (Python)-type
clustered map from \texttt{diagrams-python/fig\_09\_concern\_locations.py}. Four
\texttt{dgcluster} containers, thirteen concern badges anchored north east on
their tiles with a 1 mm outset, and the single cut edge terminated by a
\texttt{\textbackslash pxmark} cross labelled "no route exists during a
procedure".}
